\section{Verification Procedure}
\label{sec:verification}

Any party (a browser, a CLI, an in-cluster verifier, or a third-party auditor)
appraises a Trust Record with no callback to the issuer and no vendor in the trust
path beyond the silicon root and the transparency-log operator. The steps are
fail-closed and ordered so that nothing is trusted before the signature is checked:

\begin{enumerate}[leftmargin=1.5em]
  \item \textbf{Signature binding.} Verify the record's signature
    (Section~\ref{sec:model:sig}) against the key in \texttt{cnf}, \emph{before any
    other field is trusted}. A record with no verifiable binding is rejected.
  \item \textbf{Freshness.} Confirm \texttt{iat} is within the maximum age (default
    24~hours unless the profile overrides). If the verifier issued a challenge
    nonce, confirm \texttt{runtime.nonce} echoes it.
  \item \textbf{Silicon root.} Confirm the attestation signature chain resolves to a
    known silicon root (AMD~\citep{amdsevsnp2023}, Intel~\citep{inteltdx2023},
    NVIDIA, or equivalent), with unrevoked vendor keys.
  \item \textbf{Reference values.} Confirm the \texttt{runtime} measurements match
    the published Reference Integrity Manifests (RIMs).
  \item \textbf{Policy.} Confirm \texttt{policy.bundle\_hash} matches the policy the
    relying party expects, and inspect \texttt{enforcement\_mode}.
  \item \textbf{Transparency.} Confirm the \texttt{transparency} receipt resolves on
    the named append-only log.
  \item \textbf{Provenance.} Confirm \texttt{build\_provenance} resolves to a
    trusted builder at the required SLSA level.
\end{enumerate}

\subsection{Result object}

The verifier emits its own appraisal in the EAR format, separating \emph{what was
claimed} from \emph{what was accepted}. The appraisal carries an overall
\texttt{status} (\texttt{affirming}, \texttt{warning}, \texttt{contraindicated},
or \texttt{none}) together with a per-claim breakdown of which fields verified and
which did not. A field may be \emph{unverified} without the record being
\emph{invalid}: for instance, a silicon cert-chain step that needs an online vendor
service (AMD KDS for the SEV-SNP VCEK chain, Intel DCAP for the TDX quote) may be
unavailable to an offline verifier, which then reports that field as unverified and
downgrades the status rather than rejecting outright. A failed \emph{signature}
binding or a \emph{stale} record, by contrast, is fatal: the record is rejected. The
cMCP reference verifier (\texttt{cmcp\_verify}) implements exactly this object, with
status values \texttt{verified}, \texttt{partially\_verified}, and
\texttt{unverified}~\citep{cmcp2026}.

A Relying Party keys its decision on the appraisal status plus its own policy on
which fields are mandatory for its risk tolerance.
